\documentclass[11pt,a4paper]{article}

\usepackage[T1]{fontenc}
\usepackage[utf8]{inputenc}
\usepackage[spanish]{babel}
\usepackage{geometry}
\usepackage{microtype}
\usepackage{setspace}
\usepackage{booktabs}
\usepackage{tabularx}
\usepackage{enumitem}
\usepackage{hyperref}
\usepackage{xcolor}
\usepackage{titlesec}
\usepackage{fancyhdr}
\usepackage{array}
\usepackage{parskip}
\usepackage{longtable}

\geometry{margin=2.3cm, top=2.6cm, bottom=2.6cm}
\setstretch{1.15}
\setlist[itemize]{leftmargin=1.4em,itemsep=0.28em,topsep=0.4em}
\setlist[enumerate]{leftmargin=1.4em,itemsep=0.28em,topsep=0.4em}

\definecolor{upnavy}{HTML}{1F3A5F}
\definecolor{upgray}{HTML}{5A6570}
\definecolor{upsoft}{HTML}{E8EEF4}

\hypersetup{colorlinks=true,urlcolor=upnavy,linkcolor=upnavy}
\titleformat{\section}{\large\bfseries\color{upnavy}}{\thesection.}{0.6em}{}
\titleformat{\subsection}{\normalsize\bfseries\color{upnavy}}{\thesubsection.}{0.5em}{}
\titleformat{\subsubsection}{\normalsize\itshape\color{upnavy}}{\thesubsubsection.}{0.5em}{}

\pagestyle{fancy}
\fancyhf{}
\fancyhead[L]{\small\color{upgray}Metodología de extracción --- Defensoría}
\fancyhead[R]{\small\color{upgray}Universidad del Pacífico}
\fancyfoot[C]{\thepage}
\renewcommand{\headrulewidth}{0.3pt}

\newcolumntype{Y}{>{\raggedright\arraybackslash}X}

\begin{document}

\begin{titlepage}
\centering
\vspace*{1.4cm}
{\color{upnavy}\rule{\textwidth}{1.4pt}}\\[0.7em]
{\LARGE\bfseries\color{upnavy} Metodología de extracción de conflictos}\\[0.35em]
{\Large De los PDF de la Defensoría del Pueblo a un panel caso $\times$ mes}\\[0.7em]
{\color{upnavy}\rule{\textwidth}{1.4pt}}\\[1.3em]
{\large Contraste entre el CSV de la carpeta Minería (Drive)\\[0.2em]
y el extractor local \texttt{defensoria\_extract}}\\[2em]

\begin{flushleft}
\renewcommand{\arraystretch}{1.35}
\begin{tabular}{@{}l@{\hspace{1.2em}}l@{}}
\textbf{Proyecto} & Minería, conflictos sociales e indicadores territoriales \\
\textbf{Institución} & Universidad del Pacífico \\
\textbf{Fuente} & Reportes mensuales de conflictos sociales, Defensoría del Pueblo \\
\textbf{Disco Drive} & \texttt{Johao\_asistente} $\rightarrow$ carpeta Minería \\
\textbf{Código local} & \texttt{02\_codigo/defensoria\_extract} \\
\textbf{Salida local} & \texttt{01\_datos/procesados/conflictos/panel\_caso\_mes.csv} \\
\textbf{Fecha} & Agosto 2026 \\
\end{tabular}
\end{flushleft}

\vfill
\begin{flushleft}
\color{upgray}\small
Este documento describe, con el detalle del código y de los archivos que quedaron en Drive, cómo se construyó el CSV de conflictos de la carpeta Minería y cómo se volvió a extraer el contenido con el paquete local. No es un paper: es la bitácora de método, criterios y mejoras. El archivo \texttt{.tex} no se compiló a propósito; vive en \texttt{04\_docs}.
\end{flushleft}
\end{titlepage}

\tableofcontents
\newpage

\section{Qué se compara y por qué}

Hay dos productos que parecen el mismo ---una tabla de conflictos sacada de los reportes mensuales de la Defensoría--- y no lo son.

El primero está en el Shared Drive \texttt{Johao\_asistente}, carpeta Minería. Allí el equipo dejó, entre 2024 y febrero de 2025, un pipeline en notebooks (\texttt{02\_Codigo/conflictos\_defensoria/Extraccion\_2004\_2025}) y dos CSV de trabajo:

\begin{itemize}
  \item \texttt{Conflicto.csv}: 83 casos únicos, todos etiquetados como mineros, con resumen de ChatGPT, pueblos, minas y compañía.
  \item \texttt{ReporteFechaxConflicto.csv}: 3\,777 filas caso $\times$ mes, solo para esos 83 grupos, cubriendo ediciones 7 a 251 (septiembre 2004--enero 2025).
\end{itemize}

El segundo es el trabajo local de este directorio. El paquete \texttt{defensoria\_extract} lee los PDF nativos y escribe un panel de \textbf{35\,195} filas (caso $\times$ mes) sobre \textbf{269} reportes (n.$^\circ$~1--269, abril 2004--julio 2026), sin tirar los conflictos no mineros. La minería queda como bandera (\texttt{es\_socioambiental}, \texttt{menciona\_mineria}).

La diferencia no es cosmética. Cambia la unidad de observación, el universo, la forma de leer el PDF, el criterio de identidad de un caso y el criterio de ``esto es minería''. Lo que sigue reconstruye las dos metodologías con el detalle del código, no con una glosa.

\section{El problema común}

Los reportes no son una base de datos. Son PDF de maqueta cambiante (2004--2026): listas narrativas, tablas por región, fichas de dos columnas con \texttt{Tipo:}/\texttt{Caso:}/\texttt{Ubicación:}, infografías de ciclo. Un mismo conflicto reaparece mes a mes con redacción distinta. Extraer ``el CSV'' implica, en ambos pipelines, cinco decisiones:

\begin{enumerate}
  \item Qué PDF entran (cobertura y duplicados).
  \item Dónde empieza y termina el detalle de fichas (no la portada estadística).
  \item Cómo recortar una ficha (campos, columnas, saltos de página).
  \item Cómo saber que dos textos son el mismo caso (identidad).
  \item Cómo marcar minería (filtro que borra el resto vs.\ bandera que lo deja).
\end{enumerate}

Las dos metodologías responden esas cinco preguntas de forma distinta.

%----------------------------------------------------------------------
\section{Metodología Drive: de PDF recortado a CSV}

Ruta en Drive: \texttt{Minería/02\_Codigo/conflictos\_defensoria/Extraccion\_2004\_2025}. Notebooks: \texttt{Procesamiento 2004 - 2025.ipynb} ($\approx$498 celdas), \texttt{Unificación.ipynb}, \texttt{Fuzzy.ipynb}, \texttt{Casos\_Filtrados\_Actualizado.ipynb}. Inventario de PDF: \texttt{conflictos\_sociales\_pdfs.csv} (251 ediciones, n.$^\circ$~1--251, todas marcadas \texttt{Revisado=OK}).

\subsection{Paso 0. Conversión externa (iLovePDF)}

El código no lee el PDF como documento con geometría. El flujo de trabajo, visible en las rutas \texttt{C:\textbackslash Users\textbackslash Rofer\textbackslash \ldots\textbackslash ILOVEPDF\textbackslash Procesamiento 2004 - 2025}, es:

\begin{enumerate}
  \item Recortar el PDF original a un PDF ``de sección'' (solo activos, solo latentes, etc.).
  \item Pasar ese recorte por iLovePDF (o equivalente) a Excel o Word.
  \item Parsear el Excel/Word con pandas o \texttt{python-docx}.
\end{enumerate}

Eso tiene una consecuencia de método: \textbf{se pierde la geometría de página}. Las dos columnas de una ficha (descripción a la izquierda, hechos del mes a la derecha) llegan como celdas de una hoja, no como bloques con coordenadas $(x,y)$. Si el conversor mezcla columnas, el parser hereda el error.

Herramientas de recorte: \texttt{pdfplumber} para buscar frases de inicio/fin en el texto de página; \texttt{PyPDF2} (\texttt{PdfReader}/\texttt{PdfWriter}) para copiar el rango a \texttt{recortado\_*.pdf}.

\subsection{Paso 1. Recorte por frases de sección}

La función central es \texttt{recortar\_pdfs(frase\_inicio, frase\_fin, \ldots)}. Recorre páginas, guarda el índice de la primera que contiene la frase de inicio y el de la primera posterior que contiene la de fin, y escribe un PDF nuevo. Las frases se cambiaron a mano cada vez que la Defensoría renumeró el índice. El notebook de 2024--2025 deja anotadas, mes a mes, las rúbricas reales, por ejemplo:

\begin{itemize}
  \item Activos, setiembre 2024: ``V.\ DETALLE DE LOS CONFLICTOS SOCIALES ACTIVOS'', luego ``5.1 CONFLICTOS ACTIVOS DESARROLLADOS EN UN SOLO DEPARTAMENTO'', corte en ``5.3 CONFLICTOS REACTIVADOS''.
  \item Latentes: ``VI.\ DETALLE DE LOS CONFLICTOS LATENTES'' / ``6.1 CONFLICTOS QUE HAN PASADO DE ESTADO ACTIVO A LATENTE''.
  \item Inactivos: ``6.2 \ldots HAN SALIDO DEL REGISTRO PRINCIPAL POR INACTIVIDAD PROLONGADA''.
  \item Resueltos: ``1.4 CONFLICTOS RESUELTOS'' (y en octubre 2024 ya es ``8.1'' dentro de ``VIII.\ CASOS QUE SALIERON DEL REPORTE'').
\end{itemize}

Criterio: \textbf{una sección del índice = un archivo}. El estado (Activo, Latente, Resuelto, Inactivo, Fusionado, Nuevo) no se infiere de la ficha: se hereda de la carpeta a la que cayó el recorte.

Ventana de trabajo: once periodos con parser propio, porque la maqueta no era estable.

\begin{center}
\renewcommand{\arraystretch}{1.2}
\begin{tabularx}{\textwidth}{@{}l Y@{}}
\toprule
\textbf{Periodo} & \textbf{Qué se recortaba y cómo se leía} \\
\midrule
Abril 2004 & Solo activos; texto de Word (\texttt{python-docx}), más dos casos de Cajamarca añadidos a mano. \\
Mayo 2004 & Excel de tabla regional; se rellenaba departamento al ver filas \texttt{REGIÓN \ldots} \\
Jun--ago 2004 & Activos, nuevos y concluidos; párrafos con etiquetas Antecedentes / Cuestionamiento / Hechos. \\
Set--oct 2004 & Se suman inactivos (``Desactivado'' vs ``Resuelto''). \\
Nov 2004--dic 2005 & Activos, latentes y resueltos (estos últimos a veces en cuadros). \\
Ene--dic 2006 & Texto corrido; funciones distintas por estado. \\
Ene 2007--mar 2008 & Tablas + texto (reportes 35--49). \\
Abr 2008--jul 2011 & Cuatro estados; aparece \texttt{pdfplumber} de forma sistemática. \\
Ago 2011--jul 2014 & Latentes partidos en dos oleadas (hasta n.$^\circ$~108 y desde 109). \\
Ago--set 2014 & Se añaden fusionados. \\
Oct 2014--set 2024 & Maqueta Adjuntía; parser de dos columnas sobre Excel. \\
Oct 2024--ene 2025 & Mismos campos, rúbricas de índice distintas (ciclo / ``no registraron hechos''). \\
\bottomrule
\end{tabularx}
\end{center}

\subsection{Paso 2. Parseo del Excel convertido}

Una vez el recorte es una hoja, el parser típico de la etapa Adjuntía (\texttt{procesar\_reporte\_conflictos\_activos}) hace esto:

\begin{enumerate}
  \item Lee \emph{todas} las hojas (\texttt{sheet\_name=None}, sin encabezado).
  \item Tira filas cuyo texto coincide con títulos de sección o con las palabras ``Descripción'', ``Hechos del mes'', ``Estado actual'' (encabezados de tabla que el conversor dejó como celdas).
  \item Fuerza un esquema de dos columnas: izquierda = ficha, derecha = hechos. Según el periodo, o bien concatena todas las columnas menos la última, o bien se queda con la primera y la última.
  \item Acumula filas hasta ver un nuevo \texttt{Tipo} / \texttt{CASO} / nombre de departamento; entonces cierra la ficha anterior.
  \item Parte la descripción con regex de etiquetas (\texttt{Tipo}, \texttt{Caso}, \texttt{Ubicación}, actores) ---función \texttt{procesar\_parrafo}--- y, si hace falta, \texttt{combinar\_caso\_tipo} cuando el conversor pegó ``Tipo'' y ``Caso'' en la misma celda.
  \item Detecta diálogo en la columna de hechos (\texttt{procesar\_columna\_dialogo}).
\end{enumerate}

Criterio de inicio de ficha: la celda izquierda \texttt{startswith('Tipo')} o \texttt{startswith('CASO')}, o es un departamento de una lista canónica de 25 nombres. Criterio de corte de departamento: la misma lista; los casos ``más de un departamento'' se activan cuando aparece esa rúbrica, no por geometría.

En 2004--2007 el criterio es otro: filas numeradas (\texttt{1.\ \ldots}), bloques \texttt{REGIÓN X}, y secciones de párrafo (Antecedentes, Cuestionamiento, Hechos, Intervención defensorial). El campo \texttt{Caso} a menudo se \emph{reconstruye} concatenando Antecedentes + Hechos + Cuestionamiento, porque la maqueta temprana no trae la etiqueta \texttt{Caso:}.

\subsection{Paso 3. Unificación de esquemas}

\texttt{Unificación.ipynb} no vuelve a leer PDF. Carga un Excel por periodo y por estado, \textbf{renombra columnas a un esquema común} y concatena:

\begin{itemize}
  \item \texttt{activototal} = 12 bloques de activos (abril 2004 $\ldots$ enero 2025), más los ``nuevos'' de 2004 tratados como activos.
  \item \texttt{latentetotal}, \texttt{inactivototal}, \texttt{resueltototal}, \texttt{fusionadototal}.
\end{itemize}

Ejemplos del criterio de armonización (no es un diccionario único: es celda a celda):

\begin{itemize}
  \item Mayo 2004: \texttt{Autoridad cuestionada} $\to$ Actores; \texttt{Motivo} $\to$ Caso; \texttt{Fecha de inicio} $\to$ Ingreso como caso nuevo.
  \item Set--oct 2004: \texttt{Caso} se arma con \texttt{Antecedentes + Hechos + Cuestionamiento}.
  \item 2008--2011 latentes: se separa el tipo del texto con una lista cerrada de 30+ variantes (``conflicto socioambiental'', ``conflicto por asuntos de gobierno local'', \ldots).
  \item Resueltos 2004: \texttt{Forma de resolución} = Desenlace + Últimos acontecimientos.
\end{itemize}

Salida intermedia: \texttt{ConflictosTotales2024\_2025.xlsx} (todos los estados, todavía sin identidad fuzzy). Cada fila lleva \texttt{Estado} impuesto por el bloque (\texttt{activototal['Estado']='Activo'}, etc.), no leído del PDF.

\subsection{Paso 4. Identidad fuzzy}

\texttt{Fuzzy.ipynb}:

\begin{enumerate}
  \item Estandariza \texttt{Caso} a minúsculas, quita puntuación, \textbf{conserva tildes} (\texttt{NFC}), no quita stopwords.
  \item Cruza con el inventario de PDF para pegar fecha y número de edición.
  \item Agrupa con \texttt{rapidfuzz.fuzz.token\_sort\_ratio}, umbral \textbf{90}, sobre la \textbf{lista nacional} de textos (complejidad $n^2$, sin bloquear por departamento).
  \item Asigna \texttt{Grupo\_k}. Calcula primera/última fecha, meses entre extremos, meses faltantes, y si el grupo pasó a inactivo / fusionado / resuelto.
\end{enumerate}

Límite documentado por el propio equipo: Excel trunca celdas en 32\,767 caracteres, así que grupos muy largos (muchas variantes concatenadas con \texttt{;}) se cortan. Nota en \texttt{Anotaciones - Detalles a mejorar.docx}: ``Hay una limitación de 32767, para agrupar''.

\subsection{Paso 5. Filtro de minería (lista cerrada + revisión)}

\texttt{Casos\_Filtrados\_Actualizado.ipynb} no usa el tipo ``socioambiental'' como filtro principal. Limpia tildes con \texttt{unidecode} y busca \textbf{palabras exactas} (límites de palabra) de una lista armada a mano:

\begin{itemize}
  \item Unidades / compañías: Antamina, Antapaccay, Atacocha, Cerro Corona, Las Bambas, Pucamarca, Quellaveco, San Cristóbal, Toromocho, Yanacocha, Glencore, Nexa, Gold Fields, MMG, Minsur, Anglo, Volcan, Chinalco, Newmont, Xstrata.
  \item Pueblos asociados a esas minas (San Marcos, Challhuahuacho, Hualgayoc, Nueva Morococha, \ldots).
\end{itemize}

Quien matchea pasa a \texttt{FuzzyCompletar\_filtrado}. Después hay etiquetado manual (y resumen ChatGPT) de \texttt{EsMinera}, \texttt{EsIlegal}, \texttt{EsInformal}, \texttt{Pueblos}, \texttt{Minas}, \texttt{Compañía minera}. El CSV final \texttt{Conflicto.csv} tiene 83 filas, \textbf{todas} \texttt{EsMinera = Sí}.

El panel \texttt{ReporteFechaxConflicto.csv} no es el universo de conflictos de la Defensoría: es la serie de esos 83 grupos (3\,777 filas; 3\,208 activos, 508 latentes, 34 resueltos, 26 inactivos, 1 fusionado; 244 PDF distintos).

\subsection{Errores que el propio pipeline dejó escritos}

El Word \texttt{Anotaciones - Detalles a mejorar.docx} es parte del método: lista fallos conocidos y no corregidos al momento de cerrar el CSV.

\begin{itemize}
  \item Faltan activos de mayo 2006 (n.$^\circ$~27) ``por el inicio de la página inicial''.
  \item Faltan latentes marzo 2006--marzo 2007 (error de tabulación).
  \item Resueltos enero 2007--marzo 2008 se recortaron con la frase de latentes.
  \item 2007--2011: el match de ``Actor'' corta la palabra ``directorio''.
  \item Latentes 2007--2008 incompletos porque la sección no agarra el párrafo entero.
  \item Diciembre 2012 no corrió bien; Corire no debería compartir caso con Racrachaca.
  \item Tipos que el filtro no cubría: laboral, gobierno local (aparecen igual, residuales, en el CSV de 83).
\end{itemize}

Eso no es un detalle menor: el CSV de Drive es un \textbf{producto curado y filtrado}, con huecos conscientes, no un volcado exhaustivo de fichas.

%----------------------------------------------------------------------
\section{Metodología local: el paquete \texttt{defensoria\_extract}}

El extractor local no pasa por iLovePDF. Lee el PDF con PyMuPDF (\texttt{fitz}), clasifica la maqueta, arma fichas con geometría o con texto corrido, y escribe CSV. Entrada del CLI: \texttt{python -u 02\_codigo/extraer\_conflictos.py}.

\subsection{Paso 0. Corpus e inventario (antes de extraer fichas)}

Antes del parseo hay tres scripts que el pipeline de Drive no tenía como capa formal:

\begin{enumerate}
  \item \texttt{consolidar\_pdfs\_defensoria.py} / \texttt{completar\_faltantes\_defensoria.py}: unen la descarga web con los PDF de Drive Minería. Resultado: 275 PDF en \texttt{01\_datos/crudos/pdfs}; 269 mensuales n.$^\circ$~1--269; falta solo junio 2005. Se excluyen especiales y el archivo \texttt{MALNOMBRADO\_\ldots}.
  \item \texttt{extraer\_fechas\_reportes.py}: número, mes y año desde portada (prioridad) y nombre de archivo. Criterio: regex de ``reporte \ldots n.$^\circ$~$k$'', ``al $d$ de $mes$ de $año$'', y mes/año pegado al encabezado (formato 2025).
  \item \texttt{clasificar\_formatos\_pdf.py}: una etiqueta de maqueta por PDF, con reglas sobre las dos primeras páginas (palabras clave, si está escaneado, si menciona ``fase temprana''/``desescalamiento'', si dice Adjuntía, Sumilla, Unidad, ``distinta intensidad'').
\end{enumerate}

Seis formatos de trabajo (más un especial que el extractor salta):

\begin{center}
\renewcommand{\arraystretch}{1.15}
\begin{tabular}{@{}lrrll@{}}
\toprule
\textbf{Formato} & \textbf{PDF} & \textbf{n.$^\circ$} & \textbf{Años} & \textbf{Modo de lectura} \\
\midrule
\texttt{lista\_narrativa\_2004} & 1 & 1 & 2004 & texto corrido \\
\texttt{reporte\_n\_tabla\_regiones} & 33 & 2--34 & 2004--06 & texto corrido \\
\texttt{unidad\_reporte\_n} & 15 & 35--49 & 2007--08 & texto corrido \\
\texttt{unidad\_sumilla} & 14 & 50--62 & 2008--09 & columnas \\
\texttt{adjuntia\_portada\_estadistica} & 177 & 63--238 & 2009--23 & columnas \\
\texttt{infografia\_ciclo\_2024} & 33 & 239--269 & 2024--26 & columnas \\
\bottomrule
\end{tabular}
\end{center}

Criterio de modo: si el formato está en \{\texttt{lista\_narrativa\_2004}, \texttt{reporte\_n\_tabla\_regiones}, \texttt{unidad\_reporte\_n}\} $\to$ texto corrido. Si está en \{sumilla, adjuntía, infografía\} $\to$ columnas. Si hay duda, se cuenta \texttt{tipo:} en las primeras 55 páginas: $\ge 8$ apariciones implica columnas.

Si hay dos PDF del mismo número, se queda el de más páginas (\texttt{cargar\_inventario}).

\subsection{Paso 1. Lectura geométrica de la página (\texttt{layout.py})}

\texttt{bloques\_pagina} pide a PyMuPDF \texttt{get\_text("dict")}: bloques de texto con \texttt{bbox} $(x_0,y_0,x_1,y_1)$. Criterios de limpieza \emph{antes} de parsear fichas:

\begin{itemize}
  \item Se descartan bloques de imagen (\texttt{type $\ne$ 0}).
  \item Cabecera: $y_1 < 55$; pie: $y_0 >$ alto $-50$.
  \item Se tiran encabezados de página por contenido plegado (sin tildes): ``reporte mensual de conflictos'', ``adjuntía para la prevención'', ``subadjuntía\ldots'', ``dirección de la unidad de conflictos'', y etiquetas sueltas (\texttt{descripcion}, \texttt{hechos del mes}, \texttt{estado actual}, \texttt{n}, \texttt{departamento}).
  \item Números de página de $\le 3$ dígitos.
  \item Columna: si el punto medio $x$ del bloque es $< 300$, izquierda (\texttt{L}); si no, derecha (\texttt{R}). Orden de lectura: $y$, luego L antes que R, luego $x_0$.
\end{itemize}

Departamentos: diccionario canónico de 25 nombres + Multirregional / Multidepartamental / Nacional, con tildes unificadas (Áncash, Apurímac, \ldots) y soporte de ``AREQUIPA / PUNO''. Se acepta un bloque como departamento solo si, tras limpiar, cae en ese diccionario y el texto tiene menos de 48 caracteres.

Esto sustituye al recorte por frase + Excel: la página se lee como la Defensoría la diagramó.

\subsection{Paso 2. Ventana de detalle (\texttt{secciones.py})}

No se recorta un PDF hijo. Se recorre el documento completo y se decide, página a página, si ya se está en el detalle.

\textbf{Inicio} (\texttt{pagina\_inicia\_detalle}), sobre texto plegado. Criterios (cualquiera basta, con guardas):

\begin{itemize}
  \item ``detalle de los conflictos sociales activos'' y además hay $\ge 1$ \texttt{tipo:} o la rúbrica ``en un (solo) departamento''.
  \item ``detalle de los conflictos'' con $\ge 2$ \texttt{tipo:}.
  \item Anexo + ``descripcion de los conflictos''; o ``iii.\ descripcion de los conflictos''.
  \item ``1.1 conflictos activos'' con $\ge 2$ \texttt{tipo:}.
  \item $\ge 2$ \texttt{tipo:} y la frase ``la defensoria del pueblo da cuenta'' + ``conflictos activos''.
  \item ``1.\ conflictos activos'' y ``estado:''; o ``1.\ nuevos casos''.
  \item ``conflictos vigentes'' + ``region \ldots''; o par autoridad/ubicación; o par autoridad/motivo.
\end{itemize}

Guardas para \emph{no} abrir el detalle en la sumilla o en cuadros consolidados: si hay ``sumilla'' y $<3$ \texttt{tipo:}; si hay ``informacion consolidada'' y $<2$ \texttt{tipo:}; si coinciden ``fecha de inicio'' y ``estado situacional''.

\textbf{Fin} (\texttt{pagina\_es\_stop}): solo después de haber arrancado. Se mira el encabezado de la página (primeras 12 líneas). Frases de corte: acciones colectivas de protesta, hechos de violencia contra la vida, actuaciones defensoriales, acciones de violencia subversiva, alertas tempranas. No se corta si la página sigue teniendo \texttt{tipo:} o si ``sumilla'' aparece en el primer tramo (sumillas embebidas).

\textbf{Estado del bloque} (\texttt{estado\_de\_bloque}): títulos cortos ($<220$ caracteres) que cambian el estado de las fichas siguientes: Activo, Latente, Fusionado, Inactivo, Reactivado, Resuelto, Nuevo. ``Desarrollados en más de un departamento'' pone departamento = Multidepartamental y estado Activo.

\subsection{Paso 3a. Fichas en dos columnas (\texttt{ExtractorPaginas})}

Clase de estado. Recorre bloques L/R intercalados y mantiene una ficha abierta que \textbf{puede cruzar páginas} (el recorte de Drive cortaba el PDF, así que una ficha partida entre el final de un recorte y el inicio de otro se perdía o se partía).

Reglas:

\begin{itemize}
  \item Título de estado o nombre de departamento: actualiza el contexto, no se pega a la ficha.
  \item ``Observación'' / ``Cuadro n'' se ignoran.
  \item Columna derecha: hechos del mes (y badges: ``caso nuevo'', ``hay diálogo'', ``no hay diálogo'').
  \item Columna izquierda: cuerpo. Si el bloque trae más de un \texttt{Tipo:}, se parte. Un bloque que cumple \texttt{es\_inicio\_ficha} abre ficha nueva, salvo el caso en que la ficha abierta todavía no tiene \texttt{Tipo:} y el trozo es precisamente ese \texttt{Tipo:} (continuación, no ficha nueva).
\end{itemize}

Inicio de ficha: línea que empieza por \texttt{Tipo:} o \texttt{Código:}; ``caso nuevo'' / ``caso reactivado''; ``Caso n$^\circ$ $k$''; en modo numerado, \texttt{$k$. texto} que no sea \texttt{$k.k$} (secciones).

Al cerrar: se llama \texttt{parsear\_ficha}. Se descarta si el cuerpo es un badge, si tiene $<12$ caracteres, si el ``caso'' empieza por ``defensoría del pueblo'', o si no hay caso/tipo/código y el cuerpo es corto.

\subsection{Paso 3b. Fichas en texto corrido}

Para 2004--2008. Se parte el flujo de líneas por: títulos de estado, nombres de departamento, números de caso (\texttt{$k$.} o \texttt{$k$. texto}), o inicio de ficha. Se ignoran títulos corridos (Anexo, Descripción de los conflictos, Conflictos activos/latentes/resueltos, Nuevos casos, \ldots). El cuerpo se manda al mismo \texttt{parsear\_ficha}.

\subsection{Paso 4. Campos (\texttt{parsear\_ficha})}

Preproceso: inserta saltos de línea \emph{antes} de etiquetas que el PDF pegó en la misma línea (\texttt{tipo:}, \texttt{ubicación:}, actores primarios/secundarios/terciarios, \texttt{ingresó como}, \texttt{código}, y \texttt{caso:} cuando no va precedido de ``como ''). Criterio: las etiquetas se anclan al \textbf{inicio de línea} (\texttt{\^{}label}), para no capturar ``como caso nuevo'' como campo Caso.

Campos, en orden, con lookahead hasta la siguiente etiqueta:

\texttt{tipo}, \texttt{codigo}, \texttt{ingreso\_como}, \texttt{caso}, \texttt{ubicacion}, \texttt{actores\_primarios/secundarios/terciarios}, \texttt{actores} (si no hay primarios; o \texttt{autoridad} en maquetas 2004), \texttt{ubicacion} de \texttt{población} si faltaba, \texttt{caso} de \texttt{motivo}/\texttt{cuestionamiento} si faltaba, \texttt{hechos} (columna derecha, o campo Hechos/Antecedentes).

Si aún no hay caso: primer párrafo $>40$ caracteres que no empiece por tipo/codigo. Tipo se recorta a 160 caracteres y se corta si aparece ``ingres'' (etiqueta pegada). Numeración inicial \texttt{$k$.} se quita del caso.

Basura posterior (\texttt{\_es\_basura}): casos que empiezan por ``y fecha de inicio'', ``ultimos acontecimientos'', ``ubicacion autoridad'', ``motivo y fecha''; o filas sin caso, tipo ni código.

\subsection{Paso 5. Banderas de minería (\texttt{flags.py}), no filtro}

No se borra nada. Sobre un blob plegado (tipo + código + caso + ubicación + actores + texto crudo):

\begin{itemize}
  \item \texttt{es\_socioambiental}: el \texttt{tipo} contiene ``socioambiental'' / ``socio ambiental''.
  \item \texttt{menciona\_mineria}: palabra completa de un léxico (mina/minero/minería, concesión minera, tajo, relave, campamento/unidad/proyecto/exploración minera, hidrocarburo, petróleo, Camisea, gasoducto, Dorato, Antamina, Yanacocha, Las Bambas, Southern, Chinalco, Barrick). Si no hay match, se prueba metales como palabra completa: oro, cobre, zinc, plata ---con el comentario explícito de no marcar ``foro'' ni ``plataforma''.
\end{itemize}

Criterio de uso, escrito en el resumen de corrida: filtrar \texttt{es\_socioambiental=1} y/o \texttt{menciona\_mineria=1}; el resto queda en el panel.

Departamento vacío se completa buscando el nombre canónico (los más largos primero) en ubicación + caso.

\subsection{Paso 6. Identidad (\texttt{identidad.py})}

\begin{itemize}
  \item Estandariza: NFKD sin diacríticos, minúsculas, solo \texttt{[a-z0-9]} y espacios, quita stopwords españolas (el, la, de, del, y, en, \ldots).
  \item Agrupa \textbf{dentro de cada departamento} (no $n^2$ nacional). Umbral \texttt{token\_sort\_ratio} \textbf{88} (\texttt{rapidfuzz}; si no está, \texttt{difflib}).
  \item Textos vacíos reciben id propio.
  \item Escribe \texttt{caso\_id} y \texttt{caso\_estandarizado}.
\end{itemize}

Salidas: \texttt{panel\_caso\_mes.csv} (una fila = un caso en un mes), \texttt{casos\_unicos.csv} (un id, con $n$ meses y las banderas OR a lo largo del tiempo), \texttt{fallos.csv}, \texttt{RESUMEN.txt}.

Corrida documentada: 269 PDF, 35\,195 filas, 35\,118 con campo caso, 4\,337 ids fuzzy, 20\,104 socioambientales, 20\,173 con mención de minería, 17\,326 con ambas banderas, 0 PDF a cero fichas.

Limitación explícita del extractor (no ocultada): desde $\approx$2018 el detalle a veces lista solo activos con hechos del mes; latentes en tablas compactas sin \texttt{Tipo:} no salen como filas.

%----------------------------------------------------------------------
\section{Criterios, lado a lado}

\begin{center}
\renewcommand{\arraystretch}{1.22}
\begin{tabularx}{\textwidth}{@{}l Y Y@{}}
\toprule
\textbf{Decisión} & \textbf{Drive (CSV Minería)} & \textbf{Local (\texttt{defensoria\_extract})} \\
\midrule
Unidad de análisis & 83 casos mineros curados; panel solo de esos grupos & Todos los conflictos; minería = bandera \\
PDF & 251 (hasta n.$^\circ$~251, ene 2025) & 269 mensuales (hasta n.$^\circ$~269, jul 2026) \\
Lectura & Recorte + iLovePDF $\to$ Excel/Word & PyMuPDF, bloques con \texttt{bbox} \\
Dónde está el detalle & Frases literales del índice, reescritas por periodo & Detector de página con guardas (sumilla, consolidado) \\
Estado & Carpeta del recorte (impuesto al concatenar) & Título de sección visto en el PDF \\
Dos columnas & Primera/última celda del Excel convertido & $x$ medio menor o mayor que 300 \\
Ficha que cruza página & Frágil (el recorte parte el PDF) & Estado persistente entre páginas \\
Etiquetas & \texttt{startswith} + regex sobre celda concatenada; anclaje débil & Regex anclada a inicio de línea; lookahead a la siguiente etiqueta \\
Identidad & Fuzzy nacional, umbral 90, tildes conservadas & Fuzzy por departamento, umbral 88, sin tildes ni stopwords \\
Minería & Lista cerrada de minas/pueblos + etiqueta manual \texttt{EsMinera} & Léxico + tipo socioambiental; no se borra el resto \\
Reproducibilidad & Rutas de un PC (\texttt{C:\textbackslash Users\textbackslash Rofer\textbackslash \ldots}), 498 celdas, iLovePDF & CLI único, rutas en \texttt{rutas.py} \\
Control de calidad & Word de fallos conocidos; Excel 32\,767 & \texttt{fallos.csv}; 0 PDF vacíos; inventario de formatos \\
\bottomrule
\end{tabularx}
\end{center}

\subsection{Criterios de inclusión del CSV de Drive (los 83)}

Un grupo entra al CSV de Minería si, después del fuzzy, su texto estandarizado contiene al menos una palabra de la lista de unidades/compañías/pueblos, y luego se confirma \texttt{EsMinera=Sí}. Eso sesga el universo hacia las minas que el equipo ya tenía en el catálogo (Yanacocha, Antamina, Bambas, Toromocho, Quellaveco, \ldots). Un conflicto minero artesanal o una unidad que no está en la lista no entra, aunque el tipo sea socioambiental.

\subsection{Criterios de inclusión del panel local}

Entra toda ficha que sobreviva a basura y tenga caso, tipo o código (o cuerpo largo). El analista recorta después:

\begin{itemize}
  \item Socioambiental estricto: \texttt{es\_socioambiental=1} (20\,104 filas).
  \item Mención minera (incluye hidrocarburos del léxico): \texttt{menciona\_mineria=1} (20\,173).
  \item Intersección: 17\,326.
\end{itemize}

No hay equivalencia numérica con 3\,777: aquel número es la historia mensual de 83 grupos ya filtrados; este es el panel completo.

%----------------------------------------------------------------------
\section{Qué se mejoró y dónde está la ganancia}

\subsection{Se dejó de convertir el PDF}

La ganancia más grande no es un umbral fuzzy. Es no entregar el layout a un conversor. iLovePDF aplana dos columnas en celdas; el parser de Drive tenía que \emph{adivinar} cuál columna era hechos del mes (a veces la última, a veces ``todo menos la última''). PyMuPDF conserva $x$ y $y$: hechos van a la derecha porque están a la derecha.

Eso elimina una familia entera de errores que el Word de anotaciones ya había visto: secciones que no agarran el párrafo, tabulaciones que se comen latentes, ``Actor'' que corta ``directorio''.

\subsection{Un detector de sección en lugar de once tijeras}

Drive mantenía listas de frases de inicio/fin por periodo (y las reescribía cuando el índice pasaba de ``V.\ DETALLE\ldots'' a ``VI.\ DETALLE\ldots''). Localmente hay un detector con varias condiciones y \emph{guardas} para no abrir el detalle en la sumilla. Cuando la Defensoría renumera el capítulo, no hace falta un notebook nuevo: hace falta que alguna de las condiciones siga siendo verdad (y, de hecho, ``detalle de los conflictos sociales activos'' + presencia de \texttt{tipo:} cubre 2009--2026).

Costo consciente: las tablas compactas de latentes sin \texttt{Tipo:} no se vuelven filas. Drive sí las perseguía con recortes específicos; a cambio, fallaba cuando la frase no coincidía.

\subsection{La ficha es un objeto con estado, no una acumulación de filas de Excel}

\texttt{ExtractorPaginas} abre/cierra fichas y sobrevive al salto de página. El parser de Drive acumulaba celdas hasta ver otro \texttt{Tipo} o un departamento; si el conversor partía ``Tipo: Socioambiental'' a mitad, o si el recorte de PDF cortaba la ficha, se generaban filas rotas o se perdía el final.

El anclaje \texttt{\^{}label} evita el falso campo Caso en ``ingresó como caso nuevo'', un error clásico de regex no anclada.

\subsection{Identidad más barata y menos mezclona}

Fuzzy nacional a umbral 90 mezcla casos homónimos de distintos departamentos (Corire vs.\ Racrachaca es el ejemplo que el propio equipo anotó). Bloquear por departamento reduce comparaciones y ese tipo de colisión. Bajar a 88, quitar tildes y stopwords compensa la redacción mes a mes (``el caso de la comunidad \ldots'' vs.\ ``comunidad \ldots'') sin exigir que el texto sea casi idéntico.

No es un clustering perfecto: dos conflictos distintos en el mismo departamento con nombres parecidos pueden colapsar. El criterio es explícito y el \texttt{caso\_estandarizado} queda en el CSV para auditar.

\subsection{Minería como bandera}

El CSV de Drive es útil para un tablero de \emph{las} minas del proyecto (83 historias). Es un mal insumo si la pregunta es ``cuántos conflictos socioambientales hay en el reporte $t$''. El panel local no decide la pregunta por el usuario: deja 35\,195 filas y dos bits.

El léxico local es más abierto que la lista de unidades (cualquier ``minera'', ``relave'', ``tajo''), y más disciplinado en metales (palabra completa). Incluye hidrocarburos (Camisea, gasoducto): si el analista quiere solo metálicos, debe recortar. Drive, al filtrar por nombre de mina grande, casi no tiene ese problema ---y se pierde el resto.

\subsection{Cobertura y operación}

\begin{itemize}
  \item Dieciocho reportes más (252--269) y el tramo 2004 que el panel minero de Drive ni siquiera usaba (el \texttt{ReporteFechaxConflicto} empieza en la edición 7).
  \item Un comando, sin pasos manuales de recorte ni de conversión.
  \item Inventario de formatos reproducible; hardlinks por maqueta en \texttt{pdfs\_por\_formato}.
  \item 0 PDF mensuales a cero fichas, frente a una lista de huecos fechados en Drive.
\end{itemize}

\subsection{Lo que Drive hacía mejor (y no se copió a ciegas)}

Conviene no vender el extractor local como reemplazo total del trabajo de Minería.

\begin{itemize}
  \item Drive \textbf{curó} minas, pueblos y compañías y escribió un resumen legible (\texttt{CasoResumidoChatGPT}). El panel local no geocodifica ni nombra la unidad minera.
  \item Drive guardaba campos de tránsito que el local no prioriza: motivo de fusión, motivo de pase a latente/inactivo, forma de resolución, diálogo como columna propia (el local deja el badge en \texttt{hechos\_del\_mes}).
  \item Drive intentó extraer latentes tabulares. El local los declara fuera de alcance cuando no hay \texttt{Tipo:}.
\end{itemize}

La mejora es de \emph{extracción} (leer el PDF de forma estable y completa). El CSV de 83 sigue siendo un producto de \emph{investigación} (universo minero del grupo, etiquetas manuales). Se pueden cruzar: filtrar el panel local por \texttt{menciona\_mineria} y luego aplicar, si se quiere, el catálogo de minas de Drive.

%----------------------------------------------------------------------
\section{Números de contraste}

\begin{center}
\renewcommand{\arraystretch}{1.2}
\begin{tabular}{@{}lrr@{}}
\toprule
\textbf{Magnitud} & \textbf{Drive} & \textbf{Local} \\
\midrule
PDF en inventario de trabajo & 251 & 269 mensuales (275 en carpeta) \\
Última edición & 251 (ene 2025) & 269 (jul 2026) \\
Filas caso $\times$ mes & 3\,777 (solo 83 grupos) & 35\,195 (universo) \\
Casos únicos & 83 (mineros curados) & 4\,337 (fuzzy) \\
Activos en el panel & 3\,208 & 32\,567 \\
Latentes & 508 & 2\,213 \\
Resueltos & 34 & 402 \\
PDF distintos en el panel & 244 & 269 \\
PDF a cero fichas & varios, listados en anotaciones & 0 \\
\bottomrule
\end{tabular}
\end{center}

Fichas locales por formato: adjuntía 26\,300; infografía 2024--26 2\,893; tabla regiones 2\,604; unidad reporte 1\,743; sumilla 1\,608; lista 2004 47.

%----------------------------------------------------------------------
\section{Cómo reproducir el panel local}

No hace falta iLovePDF. Con el corpus en \texttt{01\_datos/crudos/pdfs} y el inventario de formatos:

\begin{verbatim}
python -u 02_codigo/extraer_conflictos.py
\end{verbatim}

Opciones: \texttt{--max N} (prueba), \texttt{--solo-numero 169}, \texttt{--sin-fuzzy}. Rutas en \texttt{02\_codigo/rutas.py}. Salida en \texttt{01\_datos/procesados/conflictos/}.

El paquete, por archivo:

\begin{itemize}
  \item \texttt{layout.py} --- bloques, columnas, departamentos, encabezados.
  \item \texttt{secciones.py} --- inicio/fin de detalle y estado.
  \item \texttt{fichas.py} --- parseo de campos y el extractor de páginas.
  \item \texttt{flags.py} --- socioambiental y minería.
  \item \texttt{identidad.py} --- fuzzy por departamento.
\end{itemize}

%----------------------------------------------------------------------
\section{Fuentes usadas para reconstruir el método Drive}

No se adivinó el pipeline de Minería: se leyeron los archivos vigentes en Drive.

\begin{itemize}
  \item \texttt{Procesamiento 2004 - 2025.ipynb} --- recorte, parsers por periodo.
  \item \texttt{Unificación.ipynb} --- rename/concat; escribe \texttt{ConflictosTotales2024\_2025.xlsx}.
  \item \texttt{Fuzzy.ipynb} --- umbral 90, \texttt{FuzzyCompletar.xlsx}.
  \item \texttt{Casos\_Filtrados\_Actualizado.ipynb} --- lista de minas/pueblos.
  \item \texttt{Anotaciones - Detalles a mejorar.docx} --- fallos conocidos.
  \item \texttt{conflictos\_sociales\_pdfs.csv}, \texttt{Conflicto.csv}, \texttt{ReporteFechaxConflicto.csv}.
  \item \texttt{01\_Datos/DICCIONARIO\_DATOS.md}: \texttt{conflictos\_data.json} se describe como base de conflictos mineros $\approx$ oct 2014--set 2024, cruzada con OCMAL; es otro producto, posterior al parseo de PDF.
\end{itemize}

Enlace de la carpeta Minería: \url{https://drive.google.com/drive/folders/1NfRO7wAfOVF6dasoElMsV8srW-7ThSh9}

\section{Cierre}

El CSV de la carpeta Minería es el resultado de recortar PDF por rúbricas del índice, convertirlos a Excel, parsear cada ventana 2004--2025 con un código distinto, unificar columnas a mano, agrupar textos a umbral 90 y quedarse con 83 casos que coincidían con el catálogo de minas del proyecto.

El extractor local hace la misma lectura de fondo ---fichas de la Defensoría, mes a mes--- pero sobre el PDF nativo, con geometría, un solo modelo de ficha, identidad acotada por departamento y banderas en lugar de un filtro que tira el universo. Eso es la optimización: menos conversiones, menos parsers-por-periodo, más filas auditables, y la decisión de ``qué es minería'' aplazada al análisis.

\end{document}
